\subsection{Квазиньютоновские методы}
\label{subsec:quasi-newton-anchor}

Пусть последовательность $\{x_k\}$ строится по схеме
\[
  x_k = x_{k-1} + \alpha_k p_k , \qquad k \in \NN ,
\]
причём направление спуска берётся в виде
\[
  p_k = -A_k \nabla f(x_{k-1}) = A_k\, \omega_k , \qquad k \in \NN ,
  \qquad\text{где}\qquad
  \omega_k = -\nabla f(x_{k-1}) .
\]
Здесь $A_k$ — положительно определённая матрица порядка $n$, обновляемая на
$k$-й итерации (заменитель гессиана), причём
\[
  A_k \xrightarrow[k \to \infty]{} H^{-1}(x_*) ,
  \qquad
  A_{k+1} = A_k + \Delta A_k , \quad k \in \NN ,
\]
где $\Delta A_k$ — поправочная матрица порядка $n$. Начальное значение
$A_1 = E$.

\begin{idea}
Это ровно схема $p_k = -F(x_k)\nabla f(x_k)$ из \S\ref{subsec:descent-direction},
где $F = A_k$. Идея: гессиан не вычислять вовсе, а накапливать приближение к
$H^{-1}$ по наблюдаемым парам «сдвиг аргумента — сдвиг градиента». Поправка
$\Delta A_k$ подбирается так, чтобы выполнялось \textbf{секущее уравнение}
\[
  A_{k+1}\,\Delta\omega_k = -\Delta x_k ,
\]
то есть чтобы $A_{k+1}$ действовала на уже известном направлении так же, как
действовал бы настоящий $H^{-1}$ (поскольку
$H\,\Delta x \approx \Delta \nabla f = -\Delta\omega$). Положительная
определённость $A_k$ нужна, чтобы $p_k$ оставалось направлением спуска.

Старт с $A_1 = E$ означает, что первая итерация — обычный градиентный шаг:
информации о кривизне ещё нет.
\end{idea}

Остальные методы различаются по способу построения поправочной матрицы.

\paragraph{ДФП-метод (Дэвидона--Флетчера--Пауэлла).}
\[
  A_{k+1} = A_k
    - \frac{\Delta x_k\, (\Delta x_k)\T}{\Delta\omega_k\T\, \Delta x_k}
    - \frac{A_k \Delta\omega_k\, (\Delta\omega_k)\T A_k\T}
           {\Delta\omega_k\T A_k \Delta\omega_k} .
\]

\paragraph{БФШ-метод (Бройдена--Флетчера--Голдфарба--Шенно).}
\[
  A_{k+1} = A_k
    - \frac{\Delta x_k\, (\Delta x_k)\T}{\Delta\omega_k\T\, \Delta x_k}
    - \frac{A_k \Delta\omega_k\, (\Delta\omega_k)\T A_k\T}
           {\Delta\omega_k\T A_k \Delta\omega_k}
    + \big(\Delta\omega_k\T A_k \Delta\omega_k\big)\, r^k (r^k)\T ,
\]
\[
  r^k = \frac{A_k \Delta\omega_k}{\Delta\omega_k\T A_k \Delta\omega_k}
      - \frac{\Delta x_k}{\Delta x_k\T \Delta\omega_k} .
\]

Также: \textbf{метод Пауэлла}, \textbf{метод Мак-Кормика}.

\begin{idea}
БФШ — это ДФП плюс одно слагаемое ранга единица; оба обновления сохраняют
симметрию и положительную определённость $A_k$ (при
$\Delta\omega_k\T \Delta x_k < 0$, что обеспечивается точным или достаточно
аккуратным одномерным поиском). На практике БФШ устойчивее ДФП к неточной
одномерной минимизации, поэтому именно он стал стандартным.

Знаки в формулах объясняются тем, что здесь работают с $\omega = -\nabla f$,
а не с самим градиентом: при подстановке $\Delta\omega_k = -\Delta \nabla f_k$
обе формулы превращаются в привычные записи ДФП и БФШ.
\end{idea}

\subsubsection*{Общий алгоритм}

\begin{algo}[title={Квазиньютоновский метод}]
\textbf{Предварительный этап (1):}
\begin{itemize}[nosep,left=0pt]
  \item задать $\eps$ (параметр в условии останова $\norm{\omega_k} < \eps$);
  \item выбрать начальную точку $x_0$;
  \item $A_1 = E$;
  \item $I_0 = \{n,\, 2n,\, \dots\}$ — моменты обновления алгоритма;
  \item $k := 1$;
  \item $\omega_1 := -\nabla f(x_0)$.
\end{itemize}

\medskip
\textbf{Шаг (2):}
\begin{steps}
  \item На $k$-м шаге проверка останова $\norm{\omega_k} < \eps$:\\
        если да — $x_* \approx x_{k-1}$, $f(x_*) \approx f(x_{k-1})$, конец.
  \item Вычисление направления спуска $p_k = A_k \omega_k$ и шага
        \[
          \psi_k(\alpha) = f(x_{k-1} + \alpha p_k) \to \min_\alpha ,
          \qquad
          x_k = x_{k-1} + \alpha_k p_k ,
        \]
        затем $\omega_{k+1} = -\nabla f(x_k)$.\\
        Если $k \in I_0$ — положить $A_{k+1} = E$, $k := k+1$ и вернуться
        к~шагу~1; иначе перейти к~шагу~3.
  \item $\Delta x_k := x_k - x_{k-1}$, \quad
        $\Delta\omega_k := \omega_{k+1} - \omega_k$.
  \item Вычислить $A_{k+1}$ по формуле метода, $k := k+1$, вернуться к~шагу~1.
\end{steps}
\end{algo}

\begin{idea}
Зачем сброс $A_{k+1} = E$ каждые $n$ итераций (множество $I_0$). Во-первых,
для квадратичной функции квазиньютоновские формулы восстанавливают $H^{-1}$
точно за $n$ шагов — дальше новой информации о кривизне уже не появляется.
Во-вторых, для неквадратичной $f$ накопленные поправки постепенно перестают
отвечать текущей кривизне, а ошибки округления могут нарушить положительную
определённость $A_k$. Сброс на единичную матрицу гарантированно возвращает
корректное направление спуска ценой одного градиентного шага.
\end{idea}
